\documentclass[11pt]{article}

\usepackage[margin=1in]{geometry}
\usepackage{amsmath, amssymb, amsthm, bm, mathtools}
\usepackage{booktabs, array, longtable}
\usepackage{enumitem}
\usepackage{hyperref}

\newcommand{\bY}{\mathbf{Y}}
\newcommand{\bX}{\mathbf{X}}
\newcommand{\bD}{\mathbf{D}}
\newcommand{\bF}{\mathbf{F}}
\newcommand{\bQ}{\mathbf{Q}}
\newcommand{\bB}{\mathbf{B}}
\newcommand{\bR}{\mathbf{R}}
\newcommand{\bG}{\mathbf{G}}
\newcommand{\bM}{\mathbf{M}}
\newcommand{\bL}{\mathbf{L}}
\newcommand{\bS}{\mathbf{S}}
\newcommand{\bH}{\mathbf{H}}
\newcommand{\bP}{\mathbf{P}}
\newcommand{\bA}{\mathbf{A}}
\newcommand{\bz}{\mathbf{z}}
\newcommand{\br}{\mathbf{r}}
\newcommand{\bzero}{\mathbf{0}}
\newcommand{\diag}{\operatorname{diag}}
\newcommand{\tr}{\operatorname{tr}}
\newcommand{\cov}{\operatorname{cov}}
\newcommand{\var}{\operatorname{var}}
\newcommand{\I}{\mathbf{I}}
\newcommand{\E}{\operatorname{E}}

\title{Technical Report: How the Closed-Form Update for \texorpdfstring{$\bM$}{M} Is Embedded in the Current Prop Skew-\texorpdfstring{$t$}{t} Backend}
\author{}
\date{\today}

\begin{document}

\maketitle

\section{Purpose of This Report}

This report is written as a \emph{verification memo}, not as a proposal note.
Its goal is to answer two narrow technical questions for the current prop
implementation in \texttt{D:/Github/spatial-skew-t}:

\begin{enumerate}[label=\arabic*.]
\item Is the current implementation of the closed-form maximum likelihood
estimator of the low-rank covariance matrix $\bM$ algebraically consistent with
the Tzeng--Huang theorem after the basis construction used by this backend?
\item How is that closed-form covariance update embedded inside the current
skew-$t$ sampler without replacing the rest of the Morris-style hierarchical
structure?
\end{enumerate}

The report is therefore designed to clarify what is exact, what is only a
working/profile step, and what evidence would be needed to claim that the
current code writes the intended estimator correctly.

\section{Files That Matter}

The relevant code paths are:

\begin{itemize}
\item \texttt{code/R/prop/prop\_basis.R}: basis construction, trimming, and
orthonormalization.
\item \texttt{code/R/prop/prop\_covariance.R}: closed-form update for the
low-rank block and the construction of $\bG$ and $\bR$.
\item \texttt{code/R/prop/prop\_modules.R}: high-level wrappers such as
\texttt{update\_M\_closed\_form(...)} and
\texttt{build\_R\_from\_G(...)}.
\item \texttt{code/R/prop/mcmc\_prop.R}: the hybrid skew-$t$ sampler that calls
the covariance update.
\item \texttt{code/analysis/simstudy\_prop/run-prop.R}: runner that passes
\texttt{prop\_k} and \texttt{prop\_cov\_update\_every} into the backend and
saves the resulting fit objects.
\end{itemize}

The \emph{runner} governs how jobs are launched and saved. The actual
closed-form update for $\bM$ lives in the \emph{backend}. The scientific
correctness question therefore belongs primarily to the backend, with the
runner only responsible for exposing the correct control parameters and
recording metadata.

\section{Implemented Prop Skew-\texorpdfstring{$t$}{t} Model}

\subsection{Observation-Level Form}

For replicate $t = 1,\ldots,T$, let $\bY_t \in \mathbb{R}^n$ denote the
observed response vector at the $n$ observed spatial sites. The current prop
backend can be written as
\begin{equation}
\bY_t = \bX_t \bm{\beta} + \lambda \bz_{g,t} + \bD_t \mathbf{v}_t,
\label{eq:obs-model}
\end{equation}
where:

\begin{itemize}
\item $\bX_t$ is the covariate matrix at replicate $t$;
\item $\bm{\beta}$ is the regression coefficient vector;
\item $\lambda$ is the skewness coefficient;
\item $\bz_{g,t}$ is the site-level skew latent vector induced by the current
partition/knots structure;
\item $\bD_t = \diag\{\tau_{g,t}(s_1)^{-1/2}, \ldots,
\tau_{g,t}(s_n)^{-1/2}\}$ is the scale matrix built from the current local
precisions;
\item $\mathbf{v}_t$ is the latent Gaussian spatial effect with correlation
matrix $\bR$.
\end{itemize}

Equation \eqref{eq:obs-model} matches the current code structure in
\texttt{mcmc\_prop.R}: the mean part is formed as
\texttt{x.beta + lambda * zg}, while the scale is carried by \texttt{taug}
through the standardization step.

\subsection{What Is Preserved From the Morris Hierarchy}

The current prop implementation does \emph{not} replace the whole Morris-style
hierarchy. It keeps the main non-Gaussian latent blocks:

\begin{itemize}
\item $\bm{\beta}$ update;
\item $\lambda$ update when skewness is enabled;
\item local scale update through $\tau$ and the site-level map \texttt{taug};
\item latent skew process update through \texttt{z} and \texttt{zg};
\item knot/partition updates when \texttt{nknots > 1}.
\end{itemize}

What changes is only the latent spatial dependence block $\mathbf{v}_t$. The
Mat\'ern-driven covariance update is replaced by a low-rank spatial random
effects block that is refreshed from a closed-form estimator of $\bM$ and
$\sigma_\xi^2$.

\section{Conditional Gaussian Working Residuals}

\subsection{Standardized Residuals Used by the Covariance Update}

Given the current state
\[
(\bm{\beta}, \lambda, \bz_{g,t}, \tau_{g,t}), \qquad t=1,\ldots,T,
\]
define the standardized residual vector
\begin{equation}
\br_t
=
\bD_t^{-1}
\bigl(
\bY_t - \bX_t \bm{\beta} - \lambda \bz_{g,t}
\bigr),
\qquad
t=1,\ldots,T,
\label{eq:std-resid}
\end{equation}
where $\bD_t^{-1} = \diag\{\tau_{g,t}(s_1)^{1/2}, \ldots,
\tau_{g,t}(s_n)^{1/2}\}$.

This matches the code exactly:
\[
\texttt{residuals\_std}
=
\sqrt{\texttt{taug}} \odot
\bigl(
\texttt{y} - \texttt{x\_beta} - \lambda \texttt{zg}
\bigr),
\]
implemented in \texttt{update\_residuals(...)}.

\subsection{Why Theorem 2 Is Applied to \texorpdfstring{$\br_t$}{r\_t}, Not Directly to \texorpdfstring{$\bY_t$}{Y\_t}}

Tzeng--Huang's theorem is derived for Gaussian spatial random effects
replicates with a common covariance model. In the current backend, the original
observations are not Gaussian because of the skew-$t$ hierarchy. The theorem is
therefore applied to the \emph{conditionally Gaussian working residuals}
\(\br_t\), not to \(\bY_t\) directly.

This is the first crucial interpretive point:

\begin{quote}
the current backend does not claim that the full skew-$t$ model admits a closed
form MLE for $\bM$; it claims that, conditional on the current latent blocks,
the Gaussian spatial dependence layer can be refreshed by a theorem-based
closed-form update.
\end{quote}

That is why the method should be described as a \emph{hybrid
Bayesian--profile} algorithm rather than a full Bayesian sampler for
$\bM$.

\section{Basis Construction and Why the Theorem Simplifies}

\subsection{Fixed Basis and Rank Reduction}

The current backend fixes the basis matrix through
\texttt{build\_prop\_basis(...)}. Let the requested basis size be
\(K_{\mathrm{req}}\), but let the numerically retained rank after trimming and
orthonormalization be \(r\). The code returns
\[
\bQ \in \mathbb{R}^{n \times r},
\qquad
\bQ' \bQ = \I_r.
\]

This orthonormalization step matters because it changes how Theorem~2 should be
read in the implementation. The theorem is written for a general basis matrix
\(\bF\), but in the current backend the working basis has already been
orthonormalized. Therefore the factor
\[
(\bF' \bF)^{-1/2}
\]
collapses to \(\I_r\) after basis preprocessing.

\subsection{The Working Empirical Covariance}

Collect the residual replicates as
\[
\mathbf{R}
=
\begin{bmatrix}
\br_1 & \cdots & \br_T
\end{bmatrix},
\qquad
\bS = \frac{1}{T} \mathbf{R} \mathbf{R}'.
\]
Then with the orthonormalized basis \(\bQ\), the theorem's matrix
\[
\bH
=
(\bF'\bF)^{-1/2}\bF' \bS \bF (\bF'\bF)^{-1/2}
\]
reduces to
\begin{equation}
\bH = \bQ' \bS \bQ.
\label{eq:h-simplified}
\end{equation}

The code computes exactly this quantity by
\[
\texttt{projected} = \bQ' \mathbf{R},
\qquad
\texttt{h\_mat} = \frac{1}{T}\texttt{projected}\,\texttt{projected}'.
\]
Hence
\[
\texttt{h\_mat} = \bQ' \bS \bQ.
\]

This is the second crucial point of the report:

\begin{quote}
the implementation is not omitting the \((\bF'\bF)^{-1/2}\) factor by mistake;
it has already absorbed that factor by orthonormalizing the basis before the
closed-form update is called.
\end{quote}

\section{Theorem 2 Update and Its Code Mapping}

\subsection{The Working Theorem}

Given the orthonormalized basis \(\bQ\), let
\[
\bH = \bQ' \bS \bQ = \bP \diag(d_1,\ldots,d_r)\bP'
\]
be its eigendecomposition, where \(d_1 \ge \cdots \ge d_r \ge 0\). Let
\(d_0 = 0\), and let
\[
L^* = \max \left\{
L \in \{1,\ldots,r\} :
d_L >
\max\left(
\frac{1}{n-L}\Bigl(\tr(\bS) - \sum_{k=1}^L d_k\Bigr),
\sigma_\epsilon^2
\right)
\right\},
\]
with \(L^* = 0\) if the set is empty.

Then the theorem-based update is
\begin{equation}
\widehat{\sigma}_\xi^2
=
\max\left\{
\frac{1}{n-L^*}
\left(
\tr(\bS) - \sum_{k=1}^{L^*} d_k
\right)
- \sigma_\epsilon^2,
\, 0
\right\},
\label{eq:sig-update}
\end{equation}
and
\begin{equation}
\widehat{\bM}
=
\bP \diag(\widehat d_1,\ldots,\widehat d_r)\bP',
\qquad
\widehat d_k = \max(d_k - \widehat{\sigma}_\xi^2 - \sigma_\epsilon^2, 0).
\label{eq:m-update}
\end{equation}

In the current backend, \(\sigma_\epsilon^2\) is passed as
\texttt{prop\_sigma\_eps\_sq} and defaults to \(0\).

\subsection{Implementation-Level Regularization}

The code introduces one additional regularization step beyond the theorem:
\[
\widehat{\sigma}_\xi^2 \leftarrow
\max\{\widehat{\sigma}_\xi^2, \texttt{sigma\_floor}\}.
\]
This does \emph{not} change the structure of the theorem-based update, but it
does mean the implementation is a \emph{numerically regularized} version of the
closed-form estimator rather than the exact raw formula with no lower bound.

\subsection{Theorem-to-Code Audit Table}

\begin{center}
\small
\begin{longtable}{p{0.18\textwidth} p{0.22\textwidth} p{0.23\textwidth} p{0.25\textwidth}}
\toprule
Mathematical object & Backend object & Function & Comment \\
\midrule
\endhead
\(\bQ\) & \texttt{F\_obs} & \texttt{build\_basis\_matrix(...)} & Orthonormalized observation basis \\
\(\bQ_*\) & \texttt{F\_pred} & \texttt{build\_basis\_matrix(...)} & Prediction basis rows \\
\(\br_t\) & \texttt{residuals\_std[, t]} & \texttt{update\_residuals(...)} & Standardized working residual replicate \\
\(\bS\) & implicit & \texttt{prop\_closed\_form\_update(...)} & Never formed explicitly as \(n\times n\) \\
\(\bH\) & \texttt{h\_mat} & \texttt{prop\_closed\_form\_update(...)} & Computed as \(Q'SQ\) \\
\(\bP\) & \texttt{rotation} & \texttt{prop\_closed\_form\_update(...)} & Eigenvectors of \(\bH\) \\
\((d_1,\ldots,d_r)\) & \texttt{projected\_eigs} & \texttt{prop\_closed\_form\_update(...)} & Eigenvalues of \(\bH\) \\
\(L^*\) & \texttt{lstar} & \texttt{prop\_select\_lstar(...)} & The theorem's retained index \\
\(\widehat{\sigma}_\xi^2\) & \texttt{sigma\_xi\_sq} & \texttt{prop\_closed\_form\_update(...)} & Nugget/fine-scale variance update \\
\((\widehat d_1,\ldots,\widehat d_r)\) & \texttt{lowrank\_eigs} & \texttt{prop\_closed\_form\_update(...)} & Shrunk eigenvalues \\
\(\widehat{\bM}\) & \texttt{M\_hat} & \texttt{update\_M\_closed\_form(...)} & Reconstructed as \(P \diag(\widehat d) P'\) \\
\bottomrule
\end{longtable}
\end{center}

\section{From \texorpdfstring{$\widehat{\bM}$}{M-hat} to the Working Correlation Matrix}

\subsection{Rotated-Basis Representation}

Although \texttt{update\_M\_closed\_form(...)} reconstructs
\(\widehat{\bM}\) explicitly, the subsequent covariance construction uses the
equivalent rotated-basis form. Let
\[
\bB = \bQ \bP,
\qquad
\bm{\Lambda} = \diag(\widehat d_1,\ldots,\widehat d_r).
\]
Then
\begin{equation}
\bQ \widehat{\bM} \bQ'
=
\bQ \bP \bm{\Lambda} \bP' \bQ'
=
\bB \bm{\Lambda} \bB'.
\label{eq:rot-basis}
\end{equation}

The current code uses \(\bB\) and \(\bm{\Lambda}\) rather than multiplying by
\(\widehat{\bM}\) directly. This is mathematically equivalent, and it is the
reason why \texttt{build\_R\_from\_G(...)} only requires
\texttt{rotation}, \texttt{lowrank\_eigs}, and \texttt{sigma\_xi\_sq}.

\subsection{From \texorpdfstring{$\bG$}{G} to \texorpdfstring{$\bR$}{R}}

The working covariance matrix is
\begin{equation}
\bG = \bB \bm{\Lambda} \bB' + \widehat{\sigma}_\xi^2 \I_n.
\label{eq:g-matrix}
\end{equation}
The current backend then defines
\[
\bA = \diag(\bG),
\qquad
\bR = \bA^{-1/2} \bG \bA^{-1/2}.
\]
This is implemented in \texttt{prop\_make\_cov\_state(...)}.

This standardization step is scientifically important:

\begin{itemize}
\item the theorem gives an update for \(\bG\), a covariance matrix;
\item the skew-$t$ sampler is run using \(\bR\), a correlation matrix;
\item the scale of the observations remains controlled by \(\tau\), not by
\(\bG\).
\end{itemize}

Therefore the current backend does \emph{not} optimize the full skew-$t$
likelihood over a covariance parameterization. Instead it performs a theorem-based
refresh of \(\bG\) and then converts that covariance block to a correlation
matrix before re-entering the skew-$t$ hierarchy.

\section{How the Update Is Embedded Inside the Sampler}

\subsection{Backend-Level Flow}

The runner \texttt{run-prop.R} passes two key controls into
\texttt{mcmc(...)}:

\begin{itemize}
\item \texttt{prop\_k}: requested basis size;
\item \texttt{prop\_cov\_update\_every}: how often the closed-form covariance
refresh is performed.
\end{itemize}

Inside \texttt{mcmc\_prop.R}, the relevant steps are:

\begin{enumerate}[label=\arabic*.]
\item Build the fixed basis \(\bQ\) and \(\bQ_*\).
\item Initialize the latent skew and scale blocks.
\item Compute the current working residual block
\(\{\br_t\}_{t=1}^T\).
\item Update \((\widehat{\bM}, \widehat{\sigma}_\xi^2)\) from those residuals.
\item Build \(\bR\), its inverse, log-determinant pieces, and the prediction
state.
\item Continue the skew-$t$ updates for \(\bm{\beta}\), \(\lambda\), \(\tau\),
\(\bz\), and knots.
\item Every \texttt{prop\_cov\_update\_every} iterations, refresh the
covariance block again from the new residuals.
\end{enumerate}

\subsection{Iteration-Level Pseudo-Code}

\noindent\textbf{Pseudo-code for one MCMC iteration}

\begin{enumerate}[label=\arabic*.]
\item Given current \((\bm{\beta}, \lambda, \tau, \bz, \text{knots})\), form
\(\br_t\) by \eqref{eq:std-resid}.
\item If thresholding or missingness is present, impute latent \(y\) conditional
on the current covariance state.
\item Update \(\bm{\beta}\) and \(\lambda\).
\item Update knots/partition state if knots are free.
\item Update \(\tau\).
\item Update \(\bz\) if skewness is enabled.
\item Recompute \(\br_t\).
\item If the iteration index is a covariance-refresh step, compute
\((\widehat{\bM}, \widehat{\sigma}_\xi^2)\), then rebuild \(\bR\).
\item Save the current draw, including prediction draws generated from the
updated covariance state.
\end{enumerate}

This pseudo-code is the correct interpretation of the current backend:
the closed-form estimator is embedded as a \emph{refresh step within a larger
skew-$t$ sampler}, not as a stand-alone Gaussian fitting procedure.

\section{What Is Exact, and What Is Approximate}

\subsection{Exact Statements}

The following claims are exact for the current code:

\begin{enumerate}[label=\arabic*.]
\item The basis is orthonormalized before the theorem-based update is called.
\item The matrix fed into the eigendecomposition is exactly \(Q'SQ\) for the
working residual matrix.
\item The reconstructed matrix
\(\widehat{\bM} = \bP \diag(\widehat d)\bP'\) is the theorem's closed-form
estimator in the orthonormalized basis.
\item The covariance state is built from the rotated-basis representation
\(\bB \bm{\Lambda} \bB' + \widehat{\sigma}_\xi^2 \I\), which is algebraically
equivalent to \(Q \widehat M Q' + \widehat{\sigma}_\xi^2 \I\).
\end{enumerate}

\subsection{Approximate or Profile Statements}

The following claims should \emph{not} be overstated:

\begin{enumerate}[label=\arabic*.]
\item The theorem is not being applied to the full skew-$t$ data model; it is
being applied to conditionally standardized residuals.
\item The standardization \(\bG \mapsto \bR\) is an additional modeling choice
made to preserve the scale/dependence separation.
\item The sampler does not draw \(\bM\) from a posterior distribution. It plugs
in a theorem-based update inside the MCMC loop.
\item The lower bound \texttt{sigma\_floor} is an implementation regularizer and
is not part of the raw closed-form theorem.
\end{enumerate}

Hence the right label is:

\begin{quote}
\emph{hybrid Bayesian--profile skew-$t$ inference with a theorem-based
closed-form covariance refresh}.
\end{quote}

\section{Why This Still Preserves the Skew-\texorpdfstring{$t$}{t} Process Structure}

The key scientific concern is whether introducing the closed-form covariance
refresh destroys the integrity of the skew-$t$ process. The answer is no, for
the following reason:

\begin{itemize}
\item the non-Gaussian features of the model still come from the latent skew
process, the local scale process, and the threshold/missing-data machinery;
\item the covariance refresh only replaces the Gaussian spatial dependence layer
that sits underneath those blocks;
\item prediction still uses the same skew and scale ingredients through
\texttt{lambda}, \texttt{z\_pred}, and \texttt{sigma\_pred}.
\end{itemize}

So the present backend is not ``Gaussianized'' overall. It remains a skew-$t$
process model with a low-rank Gaussian dependence component that is refreshed by
a closed-form estimator.

\section{What Must Be Verified to Claim the Estimator Is Correct}

To move from ``the code matches the intended algebra'' to ``the estimator is
implemented correctly,'' the following checks should be documented.

\subsection{Algebraic Invariants}

For each covariance refresh:

\begin{enumerate}[label=\arabic*.]
\item Verify \(Q'Q = I_r\) numerically.
\item Verify \texttt{rotation} is orthonormal.
\item Verify all \texttt{lowrank\_eigs} are nonnegative.
\item Verify \(\widehat{\bM}\) is positive semidefinite.
\item Verify \(\texttt{effective\_rank} \le \min(r, T)\).
\item Verify \(\widehat{\sigma}_\xi^2 \ge \texttt{sigma\_floor} > 0\).
\end{enumerate}

\subsection{Small-\texorpdfstring{$n$}{n}, Small-\texorpdfstring{$T$}{T} Numerical Check}

Construct a small Gaussian test case where:

\begin{itemize}
\item \(Q\) is fixed and orthonormal;
\item \(\sigma_\epsilon^2\) is known;
\item the residual replicates are truly Gaussian.
\end{itemize}

Then compare:

\begin{enumerate}[label=\arabic*.]
\item the code's \((\widehat{\bM}, \widehat{\sigma}_\xi^2)\);
\item a direct implementation of the theorem using matrix algebra outside the
backend;
\item a brute-force likelihood optimization over \(\bM\) and
\(\sigma_\xi^2\) in the same reduced-rank Gaussian model.
\end{enumerate}

If those three agree up to numerical tolerance, then the core MLE update is
very likely coded correctly.

\subsection{Sampler-Integration Check}

Independently of the theorem itself, one should also verify:

\begin{enumerate}[label=\arabic*.]
\item when \texttt{prop\_cov\_update\_every = 1}, the covariance block refreshes
at every iteration;
\item when \texttt{prop\_cov\_update\_every > 1}, the covariance state stays
fixed between refreshes;
\item the saved \texttt{fit.1\$prop} metadata is consistent with the basis rank,
\(\widehat{\sigma}_\xi^2\), and the requested \(k\).
\end{enumerate}

\section{Bottom-Line Conclusion}

From a mathematical and code-audit perspective, the current prop backend is
best described as follows:

\begin{enumerate}[label=\arabic*.]
\item The theorem-based update for \(\bM\) is implemented on the
orthonormalized basis, so the code correctly uses \(Q'SQ\) rather than the
full \((F'F)^{-1/2}F'SF(F'F)^{-1/2}\) expression.
\item The code reconstructs \(\widehat{\bM}\) in the theorem's eigenbasis and
then uses the equivalent rotated-basis form to build the covariance matrix.
\item The backend does not claim an exact skew-$t$ MLE for \(\bM\); it performs
a closed-form Gaussian covariance refresh inside a larger skew-$t$ sampler.
\item Therefore the current implementation is scientifically coherent as a
hybrid Bayesian--profile algorithm, provided the remaining numerical checks are
passed.
\end{enumerate}

The most important practical implication is this:

\begin{quote}
if the estimator is questioned, the first place to look is not the sampler
logic but the theorem-to-code mapping for the orthonormalized basis and the
\(\bG \mapsto \bR\) standardization step.
\end{quote}

\appendix

\section{Numerical Verification Appendix}

This appendix turns the previous discussion into a concrete verification
protocol. Its purpose is not to demonstrate predictive performance, but to
determine whether the current implementation of the closed-form update for
\(\bM\) is numerically consistent with the intended theorem and whether that
update is embedded correctly in the current skew-$t$ backend.

The verification logic should proceed in three layers:

\begin{enumerate}[label=\arabic*.]
\item verify the theorem-based update in a pure Gaussian reduced-rank setting;
\item verify the covariance-to-correlation conversion and prediction-state
construction;
\item verify that the closed-form step is refreshed at the intended points of
the skew-$t$ sampler.
\end{enumerate}

The order matters. If the first layer fails, then discrepancies observed in the
full skew-$t$ sampler cannot be interpreted, because the covariance update
itself has not yet been certified.

\subsection{Appendix A: Deterministic Theorem Replication Test}

\subsubsection{Goal}

The first check should isolate the theorem from the rest of the model. The
question is:

\begin{quote}
if we feed the backend an orthonormal basis and truly Gaussian replicates, does
the code return the same \((\widehat{\bM}, \widehat{\sigma}_\xi^2)\) as a direct
matrix-algebra implementation of the theorem?
\end{quote}

\subsubsection{Recommended Setup}

Use a deliberately small configuration, for example
\[
n = 12, \qquad r = 4, \qquad T = 40.
\]
Construct:

\begin{enumerate}[label=\arabic*.]
\item an orthonormal matrix \(\bQ \in \mathbb{R}^{n \times r}\), obtained for
example from a QR decomposition of a random Gaussian matrix;
\item a positive semidefinite target covariance
\[
\bM_{\mathrm{true}} = \bP_0 \diag(m_1,\ldots,m_r)\bP_0',
\]
with \(m_k \ge 0\);
\item a positive nugget/fine-scale variance \(\sigma_{\xi,\mathrm{true}}^2 > 0\);
\item the Gaussian covariance
\[
\bG_{\mathrm{true}} = \bQ \bM_{\mathrm{true}} \bQ' +
\sigma_{\xi,\mathrm{true}}^2 \I_n.
\]
\end{enumerate}

Then simulate
\[
\br_t \overset{\mathrm{iid}}{\sim} N_n(\bzero, \bG_{\mathrm{true}}),
\qquad
t=1,\ldots,T.
\]

\subsubsection{Three Parallel Calculations}

For the same simulated residual matrix \(\mathbf{R}\), compute
\((\widehat{\bM}, \widehat{\sigma}_\xi^2)\) in three ways.

\paragraph{Path 1: Backend path}
Use the current code:
\begin{enumerate}[label=\arabic*.]
\item call \texttt{update\_M\_closed\_form(F\_obs = Q, residuals\_std = R)};
\item extract \texttt{rotation}, \texttt{lowrank\_eigs},
\texttt{sigma\_xi\_sq}, and \texttt{M\_hat}.
\end{enumerate}

\paragraph{Path 2: Direct theorem path}
Implement the theorem outside the backend:
\begin{enumerate}[label=\arabic*.]
\item compute
\[
\bS = \frac{1}{T}\mathbf{R}\mathbf{R}', \qquad
\bH = \bQ'\bS\bQ;
\]
\item eigendecompose \(\bH = \bP \diag(d_1,\ldots,d_r)\bP'\);
\item compute \(L^*\), \(\widehat{\sigma}_\xi^2\), and
\(\widehat d_k = \max(d_k - \widehat{\sigma}_\xi^2, 0)\);
\item reconstruct
\[
\widehat{\bM}_{\mathrm{dir}} =
\bP \diag(\widehat d_1,\ldots,\widehat d_r)\bP'.
\]
\end{enumerate}

\paragraph{Path 3: Rebuilt rotated-basis path}
Using the backend output only, reconstruct
\[
\widehat{\bG}_{\mathrm{rot}}
=
(\bQ \,\texttt{rotation})\,
\diag(\texttt{lowrank\_eigs})\,
(\bQ \,\texttt{rotation})'
+
\texttt{sigma\_xi\_sq}\,\I_n.
\]
This path checks whether the code's rotated-basis representation is numerically
equivalent to \(Q \widehat M Q' + \widehat{\sigma}_\xi^2 I\).

\subsubsection{Acceptance Criteria}

For a successful replication check, require:

\begin{enumerate}[label=\arabic*.]
\item
\[
\max_{i,j}
\left|
\widehat M_{\mathrm{backend},ij}
-
\widehat M_{\mathrm{dir},ij}
\right|
< 10^{-8};
\]
\item
\[
\left|
\widehat{\sigma}_{\xi,\mathrm{backend}}^2
-
\widehat{\sigma}_{\xi,\mathrm{dir}}^2
\right|
< 10^{-10};
\]
\item
\[
\max_{i,j}
\left|
\widehat G_{\mathrm{rot},ij}
-
\left(
\bQ \widehat{\bM}_{\mathrm{backend}} \bQ' +
\widehat{\sigma}_{\xi,\mathrm{backend}}^2 \I_n
\right)_{ij}
\right|
< 10^{-8}.
\]
\end{enumerate}

If these fail, then the implementation of the theorem itself should be treated
as unresolved, and no higher-level skew-$t$ conclusions should be drawn.

\subsection{Appendix B: Direct Likelihood Benchmark}

\subsubsection{Goal}

The deterministic theorem replication test only shows that the code matches the
matrix algebra. A second question is whether that matrix algebra actually agrees
with the Gaussian likelihood in the reduced-rank model. This check is:

\begin{quote}
does the theorem-based update match the optimizer of the Gaussian likelihood in
the same low-rank covariance family?
\end{quote}

\subsubsection{Likelihood to Optimize}

For fixed residual replicates \(\br_1,\ldots,\br_T\), define
\[
\ell(\bM, \sigma_\xi^2)
=
-\frac{T}{2}\log |\bG|
-\frac{1}{2}\sum_{t=1}^T \br_t' \bG^{-1}\br_t,
\qquad
\bG = \bQ \bM \bQ' + \sigma_\xi^2 \I_n.
\]

To keep this benchmark feasible, use only a very small rank such as \(r=2\) or
\(r=3\), and optimize over a Cholesky parameterization
\[
\bM = \bL \bL',
\qquad
\sigma_\xi^2 = \exp(\eta),
\]
with unconstrained parameters in \(\bL\) and \(\eta\).

\subsubsection{Verification Procedure}

\begin{enumerate}[label=\arabic*.]
\item Use the same Gaussian dataset as Appendix~A.
\item Compute the theorem-based estimator
\((\widehat{\bM}_{\mathrm{thm}}, \widehat{\sigma}_{\xi,\mathrm{thm}}^2)\).
\item Independently optimize \(\ell(\bM, \sigma_\xi^2)\) numerically using
\texttt{optim(...)} or an equivalent optimizer.
\item Compare the final log-likelihood values and the resulting covariance
matrices.
\end{enumerate}

\subsubsection{Acceptance Criteria}

For the Gaussian reduced-rank benchmark:

\begin{enumerate}[label=\arabic*.]
\item the theorem-based and optimized log-likelihoods should agree to within
roughly \(10^{-6}\) to \(10^{-4}\), depending on the optimizer tolerance;
\item the resulting covariance matrices should satisfy
\[
\max_{i,j} |\widehat G_{\mathrm{thm},ij} - \widehat G_{\mathrm{opt},ij}|
\]
at a correspondingly small tolerance;
\item if likelihood values disagree materially, inspect whether the difference
comes from the \(\sigma_{\xi}^2\) floor or from optimizer failure.
\end{enumerate}

This benchmark is the strongest practical confirmation that the coded
closed-form update is not only algebraically consistent, but also likelihood
consistent in the intended Gaussian subproblem.

\subsection{Appendix C: \texorpdfstring{$\bG \mapsto \bR$}{G to R} Standardization Check}

\subsubsection{Goal}

Even if \(\widehat{\bM}\) is correct, the current backend uses \(\bR\), not
\(\bG\), in the skew-$t$ updates. Therefore the standardization step must be
verified separately.

\subsubsection{Checks to Run}

Given a covariance state built by \texttt{build\_R\_from\_G(...)}:

\begin{enumerate}[label=\arabic*.]
\item verify symmetry:
\[
\max_{i,j} |R_{ij} - R_{ji}| < 10^{-10};
\]
\item verify unit diagonal:
\[
\max_i |R_{ii} - 1| < 10^{-10};
\]
\item verify positive definiteness numerically, for example by checking that the
Cholesky factorization of \(\bR\) succeeds;
\item verify that the stored \texttt{prec\_obs} is numerically the inverse of
\(\bR\):
\[
\max_{i,j} |(\texttt{prec\_obs}\,\bR - \I)_{ij}| < 10^{-8};
\]
\item verify that the cross-covariance and conditional prediction covariance
objects are internally consistent with the Schur complement formula.
\end{enumerate}

\subsubsection{Why This Appendix Matters}

This step is where the theorem leaves its original Gaussian covariance setting
and enters the current skew-$t$ backend. Therefore, if any numerical mismatch
appears here, the theorem may still be correct while the \(\bG \mapsto \bR\)
integration is not.

\subsection{Appendix D: Sampler-Integration Verification}

\subsubsection{Goal}

The final verification layer is not about the theorem itself, but about whether
the theorem-based refresh is inserted into the MCMC loop at the intended
positions.

\subsubsection{Recommended Minimal Experiments}

Use the smallest practical runs that still exercise the covariance refresh:

\begin{enumerate}[label=\arabic*.]
\item \textbf{Gaussian no-threshold test:}
set \texttt{method = "gaussian"}, \texttt{skew = FALSE}, no censoring,
\texttt{prop\_cov\_update\_every = 1}.
\item \textbf{Skew-$t$ integration test:}
set \texttt{method = "t"}, \texttt{skew = TRUE}, with the current latent blocks
enabled, and again use \texttt{prop\_cov\_update\_every = 1}.
\item \textbf{Refresh-frequency test:}
repeat the same run with \texttt{prop\_cov\_update\_every = 3}.
\end{enumerate}

\subsubsection{What to Verify}

\paragraph{Refresh timing}
Check that:
\begin{itemize}
\item with \texttt{prop\_cov\_update\_every = 1}, the saved
\texttt{sigma\_xi} and \texttt{effective\_rank} traces can change at every
iteration after burn-in;
\item with \texttt{prop\_cov\_update\_every = 3}, the covariance state remains
unchanged between refresh points.
\end{itemize}

\paragraph{Saved metadata}
Check that the saved \texttt{fit.1\$prop} object contains:
\begin{itemize}
\item \texttt{requested\_k};
\item \texttt{basis\_kept\_cols};
\item \texttt{basis\_rank\_kept};
\item \texttt{sigma\_xi};
\item \texttt{effective\_rank};
\item \texttt{cov\_update\_every};
\item \texttt{workflow = "prop-simstudy"}.
\end{itemize}

\paragraph{Prediction path}
Check that prediction draws are produced from the refreshed covariance state and
that disabling random prediction noise (\texttt{draw = FALSE} in a local test)
returns the conditional mean implied by the stored state.

\subsection{Appendix E: Suggested Decision Rules}

For practical workflow use, the numerical checks above can be summarized into
three decision levels.

\paragraph{Green light}
Proceed with scientific experiments if:
\begin{enumerate}[label=\arabic*.]
\item Appendix~A passes exactly to numerical tolerance;
\item Appendix~B shows no meaningful likelihood gap;
\item Appendix~C confirms that \(\bR\), \(\bR^{-1}\), and the prediction state
are internally consistent;
\item Appendix~D confirms that the refresh schedule and metadata are correct.
\end{enumerate}

\paragraph{Amber light}
Proceed only with caution if:
\begin{enumerate}[label=\arabic*.]
\item the theorem replication passes, but the likelihood benchmark shows small
discrepancies traceable to \(\sigma_\xi^2\) flooring or optimizer tolerance;
\item the sampler integration works, but some metadata fields are incomplete.
\end{enumerate}

\paragraph{Red light}
Do not interpret prop-vs-Morris experiments if:
\begin{enumerate}[label=\arabic*.]
\item the theorem replication fails;
\item the rebuilt rotated-basis covariance does not match
\(Q\widehat M Q' + \widehat{\sigma}_\xi^2 I\);
\item the covariance refresh occurs at the wrong iterations;
\item the \(\bG \mapsto \bR\) state is numerically inconsistent.
\end{enumerate}

\subsection{Appendix F: Recommended Implementation Order}

To avoid conflating errors, the numerical verification should be implemented in
the following order:

\begin{enumerate}[label=\arabic*.]
\item write a stand-alone Gaussian theorem-check script that touches only
\texttt{prop\_basis.R}, \texttt{prop\_covariance.R}, and
\texttt{prop\_modules.R};
\item add a direct optimization benchmark for the same tiny Gaussian examples;
\item add a covariance-state validation script for \(\bG\), \(\bR\), and the
prediction Schur complement;
\item only then run miniature \texttt{mcmc\_prop.R} smoke tests;
\item only after all of the above pass, interpret large-scale simstudy results.
\end{enumerate}

This staged order matches the statistical logic of the method: verify the
closed-form covariance block first, then verify the correlation conversion, then
verify the skew-$t$ integration.

\begin{thebibliography}{9}

\bibitem{tzenghuang2018}
ShengLi Tzeng and Hsin-Cheng Huang.
\newblock Resolution Adaptive Fixed Rank Kriging.
\newblock \emph{Technometrics}, 60(2):198--208, 2018.

\bibitem{cressiejohannesson2008}
Noel Cressie and Gardar Johannesson.
\newblock Fixed rank kriging for very large spatial data sets.
\newblock \emph{Journal of the Royal Statistical Society: Series B},
70(1):209--226, 2008.

\end{thebibliography}

\end{document}
